% Counterexample-Guided Repair Detail — Ch. 4
% Fragment: % Counterexample-Guided Repair Detail — Ch. 4
% Fragment: % Counterexample-Guided Repair Detail — Ch. 4
% Fragment: % Counterexample-Guided Repair Detail — Ch. 4
% Fragment: \input{diagrams/tikz/repair_feedback} inside a figure environment.
% Required TikZ libraries (add to document preamble):
%   \usetikzlibrary{arrows.meta,shadows,shapes.geometric}
\usetikzlibrary{arrows.meta,shadows,shapes.geometric}

\begin{tikzpicture}[
  font=\small\sffamily,
  >=Stealth,
  vbox/.style={
    draw=#1!60!black, fill=#1!10, rounded corners=4pt,
    minimum width=6.2cm, minimum height=1.1cm,
    text width=5.8cm, align=center, inner sep=5pt,
    drop shadow={shadow xshift=0.4ex, shadow yshift=-0.4ex,
                 shadow opacity=0.30}
  },
  arr/.style={-latex, thick}
]

%----------------------------------------------------------------
% Main vertical chain
%----------------------------------------------------------------
\node[vbox=gray]   (prev) at (0, 0) {%
  \textbf{Previous Attempt}\quad
  $c_{k-1}$,\enspace $\mathit{feedback}_{k-1}$};

\node[vbox=blue]   (cand) at (0, -2.1) {%
  \textbf{New Candidate}\\[2pt]
  $c_k = \mathrm{LLM}(\mathcal{S}_t,\;\mathit{feedback}_{k-1})$};

\node[vbox=green]  (test) at (0, -4.2) {%
  \textbf{Test Execution}\\[2pt]
  Run $c_k$ on $\mathcal{D}(t) = \{x_1,\ldots,x_n\}$};

\node[vbox=red]    (cexbox) at (0, -6.7) {%
  \textbf{Counterexample}\\[2pt]
  $(x^*,\;g_t(x^*),\;f_{c_k}(x^*))$};

\node[
  draw=purple!55!black, fill=purple!8, rounded corners=4pt,
  minimum width=6.2cm, minimum height=1.3cm,
  text width=5.8cm, align=left, inner sep=5pt,
  drop shadow={shadow xshift=0.4ex, shadow yshift=-0.4ex,
               shadow opacity=0.30}
] (fbupdate) at (0, -8.9) {%
  \textbf{Feedback Update}\\[2pt]
  $\mathit{feedback}_k = \{$%
  $\mathit{counterexamples}:[\ldots],$%
  $\mathit{patterns}:[\ldots],$%
  $\mathit{attempt}:k\}$};

%----------------------------------------------------------------
% Vertical arrows with labels
%----------------------------------------------------------------
\draw[arr]
  (prev.south) --
  node[right=3pt, font=\scriptsize\sffamily] {LLM generates}
  (cand.north);

\draw[arr]
  (cand.south) --
  node[right=3pt, font=\scriptsize\sffamily] {Formal check}
  (test.north);

\draw[arr]
  (test.south) --
  node[right=3pt, font=\scriptsize\sffamily] {Some fail}
  (cexbox.north);

\draw[arr] (cexbox.south) -- (fbupdate.north);

%----------------------------------------------------------------
% RIGHT BRANCH — "All pass" → Accept check
%----------------------------------------------------------------
\node[
  draw=orange!60!black, fill=orange!8,
  diamond, aspect=2.2,
  minimum width=2.8cm, minimum height=1.0cm,
  align=center, font=\scriptsize\sffamily,
  drop shadow={shadow xshift=0.4ex, shadow yshift=-0.4ex,
               shadow opacity=0.30}
] (chk) at (5.8, -4.2) {Pattern\\clear?};

\node[
  draw=green!60!black, fill=green!10, rounded corners=4pt,
  minimum width=2.4cm, minimum height=1.0cm,
  align=center, font=\small\sffamily\bfseries,
  drop shadow={shadow xshift=0.4ex, shadow yshift=-0.4ex,
               shadow opacity=0.30}
] (ok) at (5.8, -5.9) {ACCEPTED\\$c_k$};

\draw[arr]
  (test.east) --
  node[above, font=\scriptsize\sffamily] {All pass}
  (chk.west);

\draw[arr]
  (chk.south) --
  node[right=2pt, font=\scriptsize\sffamily] {yes}
  (ok.north);

\draw[arr, dashed]
  (chk.north) -- +(0, 0.75)
  node[above, font=\scriptsize\sffamily] {no\;$\to$\;continue};

%----------------------------------------------------------------
% k > K exit arrow
%----------------------------------------------------------------
\node[
  draw=gray!55, fill=gray!5, rounded corners=3pt,
  minimum width=2.8cm, align=center, inner sep=4pt,
  font=\scriptsize\sffamily
] (exit) at (5.8, -8.9) {$k{>}K$: return best $c^*$};

\draw[arr, dashed] (fbupdate.east) -- (exit.west);

%----------------------------------------------------------------
% Curved loop arrow — Feedback Update back to Previous Attempt
%----------------------------------------------------------------
\draw[arr] (fbupdate.west) -- ++(-2.8,0) coordinate (loopW)
  (loopW) |- node[pos=0.25, anchor=east, font=\scriptsize\sffamily, align=right]
    {$k := k{+}1$} (prev.west);

\end{tikzpicture}
 inside a figure environment.
% Required TikZ libraries (add to document preamble):
%   \usetikzlibrary{arrows.meta,shadows,shapes.geometric}
\usetikzlibrary{arrows.meta,shadows,shapes.geometric}

\begin{tikzpicture}[
  font=\small\sffamily,
  >=Stealth,
  vbox/.style={
    draw=#1!60!black, fill=#1!10, rounded corners=4pt,
    minimum width=6.2cm, minimum height=1.1cm,
    text width=5.8cm, align=center, inner sep=5pt,
    drop shadow={shadow xshift=0.4ex, shadow yshift=-0.4ex,
                 shadow opacity=0.30}
  },
  arr/.style={-latex, thick}
]

%----------------------------------------------------------------
% Main vertical chain
%----------------------------------------------------------------
\node[vbox=gray]   (prev) at (0, 0) {%
  \textbf{Previous Attempt}\quad
  $c_{k-1}$,\enspace $\mathit{feedback}_{k-1}$};

\node[vbox=blue]   (cand) at (0, -2.1) {%
  \textbf{New Candidate}\\[2pt]
  $c_k = \mathrm{LLM}(\mathcal{S}_t,\;\mathit{feedback}_{k-1})$};

\node[vbox=green]  (test) at (0, -4.2) {%
  \textbf{Test Execution}\\[2pt]
  Run $c_k$ on $\mathcal{D}(t) = \{x_1,\ldots,x_n\}$};

\node[vbox=red]    (cexbox) at (0, -6.7) {%
  \textbf{Counterexample}\\[2pt]
  $(x^*,\;g_t(x^*),\;f_{c_k}(x^*))$};

\node[
  draw=purple!55!black, fill=purple!8, rounded corners=4pt,
  minimum width=6.2cm, minimum height=1.3cm,
  text width=5.8cm, align=left, inner sep=5pt,
  drop shadow={shadow xshift=0.4ex, shadow yshift=-0.4ex,
               shadow opacity=0.30}
] (fbupdate) at (0, -8.9) {%
  \textbf{Feedback Update}\\[2pt]
  $\mathit{feedback}_k = \{$%
  $\mathit{counterexamples}:[\ldots],$%
  $\mathit{patterns}:[\ldots],$%
  $\mathit{attempt}:k\}$};

%----------------------------------------------------------------
% Vertical arrows with labels
%----------------------------------------------------------------
\draw[arr]
  (prev.south) --
  node[right=3pt, font=\scriptsize\sffamily] {LLM generates}
  (cand.north);

\draw[arr]
  (cand.south) --
  node[right=3pt, font=\scriptsize\sffamily] {Formal check}
  (test.north);

\draw[arr]
  (test.south) --
  node[right=3pt, font=\scriptsize\sffamily] {Some fail}
  (cexbox.north);

\draw[arr] (cexbox.south) -- (fbupdate.north);

%----------------------------------------------------------------
% RIGHT BRANCH — "All pass" → Accept check
%----------------------------------------------------------------
\node[
  draw=orange!60!black, fill=orange!8,
  diamond, aspect=2.2,
  minimum width=2.8cm, minimum height=1.0cm,
  align=center, font=\scriptsize\sffamily,
  drop shadow={shadow xshift=0.4ex, shadow yshift=-0.4ex,
               shadow opacity=0.30}
] (chk) at (5.8, -4.2) {Pattern\\clear?};

\node[
  draw=green!60!black, fill=green!10, rounded corners=4pt,
  minimum width=2.4cm, minimum height=1.0cm,
  align=center, font=\small\sffamily\bfseries,
  drop shadow={shadow xshift=0.4ex, shadow yshift=-0.4ex,
               shadow opacity=0.30}
] (ok) at (5.8, -5.9) {ACCEPTED\\$c_k$};

\draw[arr]
  (test.east) --
  node[above, font=\scriptsize\sffamily] {All pass}
  (chk.west);

\draw[arr]
  (chk.south) --
  node[right=2pt, font=\scriptsize\sffamily] {yes}
  (ok.north);

\draw[arr, dashed]
  (chk.north) -- +(0, 0.75)
  node[above, font=\scriptsize\sffamily] {no\;$\to$\;continue};

%----------------------------------------------------------------
% k > K exit arrow
%----------------------------------------------------------------
\node[
  draw=gray!55, fill=gray!5, rounded corners=3pt,
  minimum width=2.8cm, align=center, inner sep=4pt,
  font=\scriptsize\sffamily
] (exit) at (5.8, -8.9) {$k{>}K$: return best $c^*$};

\draw[arr, dashed] (fbupdate.east) -- (exit.west);

%----------------------------------------------------------------
% Curved loop arrow — Feedback Update back to Previous Attempt
%----------------------------------------------------------------
\draw[arr] (fbupdate.west) -- ++(-2.8,0) coordinate (loopW)
  (loopW) |- node[pos=0.25, anchor=east, font=\scriptsize\sffamily, align=right]
    {$k := k{+}1$} (prev.west);

\end{tikzpicture}
 inside a figure environment.
% Required TikZ libraries (add to document preamble):
%   \usetikzlibrary{arrows.meta,shadows,shapes.geometric}
\usetikzlibrary{arrows.meta,shadows,shapes.geometric}

\begin{tikzpicture}[
  font=\small\sffamily,
  >=Stealth,
  vbox/.style={
    draw=#1!60!black, fill=#1!10, rounded corners=4pt,
    minimum width=6.2cm, minimum height=1.1cm,
    text width=5.8cm, align=center, inner sep=5pt,
    drop shadow={shadow xshift=0.4ex, shadow yshift=-0.4ex,
                 shadow opacity=0.30}
  },
  arr/.style={-latex, thick}
]

%----------------------------------------------------------------
% Main vertical chain
%----------------------------------------------------------------
\node[vbox=gray]   (prev) at (0, 0) {%
  \textbf{Previous Attempt}\quad
  $c_{k-1}$,\enspace $\mathit{feedback}_{k-1}$};

\node[vbox=blue]   (cand) at (0, -2.1) {%
  \textbf{New Candidate}\\[2pt]
  $c_k = \mathrm{LLM}(\mathcal{S}_t,\;\mathit{feedback}_{k-1})$};

\node[vbox=green]  (test) at (0, -4.2) {%
  \textbf{Test Execution}\\[2pt]
  Run $c_k$ on $\mathcal{D}(t) = \{x_1,\ldots,x_n\}$};

\node[vbox=red]    (cexbox) at (0, -6.7) {%
  \textbf{Counterexample}\\[2pt]
  $(x^*,\;g_t(x^*),\;f_{c_k}(x^*))$};

\node[
  draw=purple!55!black, fill=purple!8, rounded corners=4pt,
  minimum width=6.2cm, minimum height=1.3cm,
  text width=5.8cm, align=left, inner sep=5pt,
  drop shadow={shadow xshift=0.4ex, shadow yshift=-0.4ex,
               shadow opacity=0.30}
] (fbupdate) at (0, -8.9) {%
  \textbf{Feedback Update}\\[2pt]
  $\mathit{feedback}_k = \{$%
  $\mathit{counterexamples}:[\ldots],$%
  $\mathit{patterns}:[\ldots],$%
  $\mathit{attempt}:k\}$};

%----------------------------------------------------------------
% Vertical arrows with labels
%----------------------------------------------------------------
\draw[arr]
  (prev.south) --
  node[right=3pt, font=\scriptsize\sffamily] {LLM generates}
  (cand.north);

\draw[arr]
  (cand.south) --
  node[right=3pt, font=\scriptsize\sffamily] {Formal check}
  (test.north);

\draw[arr]
  (test.south) --
  node[right=3pt, font=\scriptsize\sffamily] {Some fail}
  (cexbox.north);

\draw[arr] (cexbox.south) -- (fbupdate.north);

%----------------------------------------------------------------
% RIGHT BRANCH — "All pass" → Accept check
%----------------------------------------------------------------
\node[
  draw=orange!60!black, fill=orange!8,
  diamond, aspect=2.2,
  minimum width=2.8cm, minimum height=1.0cm,
  align=center, font=\scriptsize\sffamily,
  drop shadow={shadow xshift=0.4ex, shadow yshift=-0.4ex,
               shadow opacity=0.30}
] (chk) at (5.8, -4.2) {Pattern\\clear?};

\node[
  draw=green!60!black, fill=green!10, rounded corners=4pt,
  minimum width=2.4cm, minimum height=1.0cm,
  align=center, font=\small\sffamily\bfseries,
  drop shadow={shadow xshift=0.4ex, shadow yshift=-0.4ex,
               shadow opacity=0.30}
] (ok) at (5.8, -5.9) {ACCEPTED\\$c_k$};

\draw[arr]
  (test.east) --
  node[above, font=\scriptsize\sffamily] {All pass}
  (chk.west);

\draw[arr]
  (chk.south) --
  node[right=2pt, font=\scriptsize\sffamily] {yes}
  (ok.north);

\draw[arr, dashed]
  (chk.north) -- +(0, 0.75)
  node[above, font=\scriptsize\sffamily] {no\;$\to$\;continue};

%----------------------------------------------------------------
% k > K exit arrow
%----------------------------------------------------------------
\node[
  draw=gray!55, fill=gray!5, rounded corners=3pt,
  minimum width=2.8cm, align=center, inner sep=4pt,
  font=\scriptsize\sffamily
] (exit) at (5.8, -8.9) {$k{>}K$: return best $c^*$};

\draw[arr, dashed] (fbupdate.east) -- (exit.west);

%----------------------------------------------------------------
% Curved loop arrow — Feedback Update back to Previous Attempt
%----------------------------------------------------------------
\draw[arr] (fbupdate.west) -- ++(-2.8,0) coordinate (loopW)
  (loopW) |- node[pos=0.25, anchor=east, font=\scriptsize\sffamily, align=right]
    {$k := k{+}1$} (prev.west);

\end{tikzpicture}
 inside a figure environment.
% Required TikZ libraries (add to document preamble):
%   \usetikzlibrary{arrows.meta,shadows,shapes.geometric}
\usetikzlibrary{arrows.meta,shadows,shapes.geometric}

\begin{tikzpicture}[
  font=\small\sffamily,
  >=Stealth,
  vbox/.style={
    draw=#1!60!black, fill=#1!10, rounded corners=4pt,
    minimum width=6.2cm, minimum height=1.1cm,
    text width=5.8cm, align=center, inner sep=5pt,
    drop shadow={shadow xshift=0.4ex, shadow yshift=-0.4ex,
                 shadow opacity=0.30}
  },
  arr/.style={-latex, thick}
]

%----------------------------------------------------------------
% Main vertical chain
%----------------------------------------------------------------
\node[vbox=gray]   (prev) at (0, 0) {%
  \textbf{Previous Attempt}\quad
  $c_{k-1}$,\enspace $\mathit{feedback}_{k-1}$};

\node[vbox=blue]   (cand) at (0, -2.1) {%
  \textbf{New Candidate}\\[2pt]
  $c_k = \mathrm{LLM}(\mathcal{S}_t,\;\mathit{feedback}_{k-1})$};

\node[vbox=green]  (test) at (0, -4.2) {%
  \textbf{Test Execution}\\[2pt]
  Run $c_k$ on $\mathcal{D}(t) = \{x_1,\ldots,x_n\}$};

\node[vbox=red]    (cexbox) at (0, -6.7) {%
  \textbf{Counterexample}\\[2pt]
  $(x^*,\;g_t(x^*),\;f_{c_k}(x^*))$};

\node[
  draw=purple!55!black, fill=purple!8, rounded corners=4pt,
  minimum width=6.2cm, minimum height=1.3cm,
  text width=5.8cm, align=left, inner sep=5pt,
  drop shadow={shadow xshift=0.4ex, shadow yshift=-0.4ex,
               shadow opacity=0.30}
] (fbupdate) at (0, -8.9) {%
  \textbf{Feedback Update}\\[2pt]
  $\mathit{feedback}_k = \{$%
  $\mathit{counterexamples}:[\ldots],$%
  $\mathit{patterns}:[\ldots],$%
  $\mathit{attempt}:k\}$};

%----------------------------------------------------------------
% Vertical arrows with labels
%----------------------------------------------------------------
\draw[arr]
  (prev.south) --
  node[right=3pt, font=\scriptsize\sffamily] {LLM generates}
  (cand.north);

\draw[arr]
  (cand.south) --
  node[right=3pt, font=\scriptsize\sffamily] {Formal check}
  (test.north);

\draw[arr]
  (test.south) --
  node[right=3pt, font=\scriptsize\sffamily] {Some fail}
  (cexbox.north);

\draw[arr] (cexbox.south) -- (fbupdate.north);

%----------------------------------------------------------------
% RIGHT BRANCH — "All pass" → Accept check
%----------------------------------------------------------------
\node[
  draw=orange!60!black, fill=orange!8,
  diamond, aspect=2.2,
  minimum width=2.8cm, minimum height=1.0cm,
  align=center, font=\scriptsize\sffamily,
  drop shadow={shadow xshift=0.4ex, shadow yshift=-0.4ex,
               shadow opacity=0.30}
] (chk) at (5.8, -4.2) {Pattern\\clear?};

\node[
  draw=green!60!black, fill=green!10, rounded corners=4pt,
  minimum width=2.4cm, minimum height=1.0cm,
  align=center, font=\small\sffamily\bfseries,
  drop shadow={shadow xshift=0.4ex, shadow yshift=-0.4ex,
               shadow opacity=0.30}
] (ok) at (5.8, -5.9) {ACCEPTED\\$c_k$};

\draw[arr]
  (test.east) --
  node[above, font=\scriptsize\sffamily] {All pass}
  (chk.west);

\draw[arr]
  (chk.south) --
  node[right=2pt, font=\scriptsize\sffamily] {yes}
  (ok.north);

\draw[arr, dashed]
  (chk.north) -- +(0, 0.75)
  node[above, font=\scriptsize\sffamily] {no\;$\to$\;continue};

%----------------------------------------------------------------
% k > K exit arrow
%----------------------------------------------------------------
\node[
  draw=gray!55, fill=gray!5, rounded corners=3pt,
  minimum width=2.8cm, align=center, inner sep=4pt,
  font=\scriptsize\sffamily
] (exit) at (5.8, -8.9) {$k{>}K$: return best $c^*$};

\draw[arr, dashed] (fbupdate.east) -- (exit.west);

%----------------------------------------------------------------
% Curved loop arrow — Feedback Update back to Previous Attempt
%----------------------------------------------------------------
\draw[arr] (fbupdate.west) -- ++(-2.8,0) coordinate (loopW)
  (loopW) |- node[pos=0.25, anchor=east, font=\scriptsize\sffamily, align=right]
    {$k := k{+}1$} (prev.west);

\end{tikzpicture}
